\documentclass[11pt,a4paper]{article}
\usepackage{amsmath,amssymb,amsthm}
\usepackage{geometry}
\usepackage{mathtools}
\usepackage{physics}
\usepackage{enumitem}
\usepackage{graphicx}
\usepackage{hyperref}

\geometry{margin=1in}

\newtheorem{theorem}{Theorem}
\newtheorem{definition}{Definition}
\newtheorem{property}{Property}

\title{\textbf{Action: Special Relativity, Electromagnetism, \\ and General Relativity}}
\author{}
\date{}

\begin{document}

\maketitle

\tableofcontents
\newpage

\section{Special Relativity}

\subsection{Lorentz Transformations}

The Lorentz transformations relate coordinates between inertial reference frames:

\begin{align}
x' &= \gamma(x - vt) \\
t' &= \gamma\left(t - \frac{vx}{c^2}\right)
\end{align}

where the Lorentz factor is:
\begin{equation}
\gamma = \frac{1}{\sqrt{1-\frac{v^2}{c^2}}}
\end{equation}

\subsection{Spacetime Interval and Metric}

The proper time interval is defined as:
\begin{equation}
c^2 d\tau^2 = -ds^2
\end{equation}

The Minkowski metric (signature convention):
\begin{equation}
\eta_{\mu\nu} = \text{diag}(-1, +1, +1, +1)
\end{equation}

\subsection{Four-Vectors}

\subsubsection{Position Four-Vector}
\begin{equation}
\bar{x} = (ct, x, y, z) = (x^0, x^1, x^2, x^3)
\end{equation}

\subsubsection{Four-Velocity}
\begin{equation}
\bar{u} = \frac{d\bar{x}}{d\tau} = \gamma(c, \vec{v})
\end{equation}

\subsubsection{Four-Momentum}
\begin{equation}
\bar{p} = m_0 \bar{u} = \left(\frac{E}{c}, \gamma m_0 \vec{v}\right) \quad \Rightarrow \quad p^0 = \frac{E}{c}
\end{equation}

\subsubsection{Four-Force}
\begin{equation}
\bar{F} = \frac{d\bar{p}}{d\tau}
\end{equation}

\section{Electromagnetism in Special Relativity}

\subsection{Four-Vectors in Electromagnetism}

\subsubsection{Four-Derivative}
\begin{equation}
\partial_\mu = \frac{\partial}{\partial x^\mu} = \left(\frac{1}{c}\frac{\partial}{\partial t}, \frac{\partial}{\partial x^i}\right)
\end{equation}

This is a Lorentz covariant derivative operator.

\subsubsection{Four-Current}
\begin{equation}
J^\mu = (c\rho, \vec{J})
\end{equation}

The continuity equation:
\begin{equation}
\partial_\mu J^\mu = 0
\end{equation}

\subsubsection{Four-Potential}
\begin{equation}
A^\mu = \left(\frac{\phi}{c}, \vec{A}\right)
\end{equation}

This transforms the same way as the position four-vector $\bar{x}$.

\subsection{Electromagnetic Field Tensor}

The electromagnetic field tensor is an antisymmetric rank-2 tensor:

\begin{equation}
F^{\mu\nu} = \begin{pmatrix}
0 & -E_x/c & -E_y/c & -E_z/c \\
E_x/c & 0 & B_z & -B_y \\
E_y/c & -B_z & 0 & B_x \\
E_z/c & B_y & -B_x & 0
\end{pmatrix}
\end{equation}

This is also called the Faraday tensor and satisfies:
\begin{equation}
F_{\mu\nu} = \partial_\mu A_\nu - \partial_\nu A_\mu
\end{equation}

\subsection{Lorentz Force}

The relativistic Lorentz force in covariant form:
\begin{equation}
f^\mu = q F^{\mu\nu} u_\nu
\end{equation}

This generalizes the familiar $\vec{F} = q(\vec{E} + \vec{v} \times \vec{B})$.

For non-relativistic velocities where $u \approx (1, \vec{v})$:
\begin{equation}
\frac{dp^i}{dt} = q E^i \quad \Rightarrow \quad F_{0i} \propto E_i
\end{equation}

\subsection{Field Transformations}

Under Lorentz transformation $F' = \Lambda F \Lambda^T$, the electric and magnetic fields transform as:

\textbf{Parallel components (unchanged):}
\begin{align}
E'_x &= E_x \\
B'_x &= B_x
\end{align}

\textbf{Perpendicular components:}
\begin{align}
E'_y &= \gamma(E_y - vB_z) & B'_y &= \gamma\left(B_y + \frac{v}{c^2}E_z\right) \\
E'_z &= \gamma(E_z + vB_y) & B'_z &= \gamma\left(B_z - \frac{v}{c^2}E_y\right)
\end{align}

\subsection{Maxwell's Equations in Covariant Form}

Maxwell's equations can be written compactly as:
\begin{equation}
\boxed{\partial_\mu F^{\mu\nu} = \mu_0 J^\nu}
\end{equation}

This encodes both Gauss's law and Ampère-Maxwell law.

\subsection{Gauge Invariance}

The electromagnetic action is invariant under gauge transformations:
\begin{equation}
A_\mu \longrightarrow A_\mu + \partial_\mu \Lambda(x)
\end{equation}

where $\Lambda(x)$ is an arbitrary scalar function.

\subsection{Action for Electromagnetism}

The unique action that is both gauge and Lorentz invariant, and produces Maxwell's equations:

\begin{equation}
S_{EM} = \int d^4x \left( -\frac{1}{4\mu_0} F_{\mu\nu}F^{\mu\nu} - A_\mu J^\mu \right)
\end{equation}

The Lagrangian density is:
\begin{equation}
\mathcal{L}_{EM} = -\frac{1}{4\mu_0} F_{\mu\nu}F^{\mu\nu} - A_\mu J^\mu
\end{equation}

\section{General Relativity Foundations}

\subsection{Motivation for Curved Spacetime}

In General Relativity, the motion of particles under gravity is determined by the metric $g_{\mu\nu}$. To find the metric for a given observer, we need field equations for $g_{\mu\nu}$.

Maxwell's theory of electromagnetism served as the guiding principle, as there is no direct Newtonian analog for the metric tensor.

\subsection{Invariant Volume Element}

\subsubsection{Derivation}

Consider a volume element transformation. We demand that the volume element be invariant:
\begin{equation}
dV = dV'
\end{equation}

In terms of the Levi-Civita symbol:
\begin{equation}
J(\bar{x}) \frac{1}{2!} \epsilon_{\alpha\beta} dx^\alpha dx^\beta = J(\bar{x}') \frac{1}{2!} \epsilon_{\mu\nu} dx'^\mu dx'^\nu
\end{equation}

where the permutation factor $(-1)^P$ accounts for cyclic ($P=0$) and anti-cyclic ($P=1$) permutations.

Under coordinate transformation $dx'^\mu = \Lambda^\mu_\alpha dx^\alpha$:
\begin{equation}
J(\bar{x}) = \det(\Lambda^\mu_\alpha)
\end{equation}

This is the Jacobian of the coordinate transformation.

\subsubsection{Connection to the Metric}

The metric transforms as:
\begin{equation}
g'_{\mu\nu} = (\Lambda^{-1})^\alpha_\mu (\Lambda^{-1})^\beta_\nu g_{\alpha\beta}
\end{equation}

Taking determinants:
\begin{equation}
\det(g') = (\det \Lambda)^{-2} \det(g)
\end{equation}

For Minkowski space $\eta_{\mu\nu}$ with $\det(\eta) = -1$:
\begin{equation}
\det(g) = (\det \Lambda)^2 \det(\eta_{\mu\nu})
\end{equation}

\subsubsection{General Volume Element}

The invariant volume element in $D$ dimensions is:
\begin{equation}
\boxed{dV = \sqrt{|g|} \, d^Dx}
\end{equation}

where $g = \det(g_{\mu\nu})$ and we use $\sqrt{-g}$ for Lorentzian signatures.

\textbf{Examples:}

\begin{itemize}
\item \textbf{2D Polar Coordinates:} $(x,y) \to (r,\theta)$
\begin{equation}
g_{\mu\nu} = \begin{pmatrix} 1 & 0 \\ 0 & r^2 \end{pmatrix}, \quad \sqrt{|g|} = r, \quad dV = r \, dr \, d\theta
\end{equation}

\item \textbf{3D Spherical Coordinates:}
\begin{equation}
ds^2 = dr^2 + r^2 d\theta^2 + r^2\sin^2\theta \, d\phi^2
\end{equation}
\begin{equation}
g_{\mu\nu} = \begin{pmatrix} 1 & 0 & 0 \\ 0 & r^2 & 0 \\ 0 & 0 & r^2\sin^2\theta \end{pmatrix}, \quad \sqrt{|g|} = r^2\sin\theta
\end{equation}
\begin{equation}
dV = r^2\sin\theta \, dr \, d\theta \, d\phi
\end{equation}
\end{itemize}

\subsection{Levi-Civita Symbol}

The totally antisymmetric Levi-Civita symbol is defined as:
\begin{align}
\epsilon_{ijk} &= +1 \quad \text{for cyclic permutations of } (1,2,3) \\
\epsilon_{ijk} &= -1 \quad \text{for anti-cyclic permutations} \\
\epsilon_{ijk} &= 0 \quad \text{if any two indices are equal}
\end{align}

\textbf{Application to determinants:}

For a $2\times2$ matrix:
\begin{equation}
\det(A) = \frac{1}{2!} \epsilon_{\mu\nu} \epsilon_{\alpha\beta} a^{\mu\alpha} a^{\nu\beta}
\end{equation}

For an $n\times n$ matrix:
\begin{equation}
\det(A) = \frac{1}{n!} \epsilon_{\mu_1 \cdots \mu_n} \epsilon_{\alpha_1 \cdots \alpha_n} a^{\mu_1\alpha_1} \cdots a^{\mu_n\alpha_n}
\end{equation}

\section{Covariant Derivatives}

\subsection{Motivation}

Consider whether $\partial_\mu g_{\alpha\beta}$ transforms as a tensor.

For a scalar field $\phi$:
\begin{equation}
\phi(x) \to \phi'(x') = \phi(x)
\end{equation}

The derivative of a scalar \emph{is} a tensor:
\begin{equation}
\partial'_\mu \phi'(x') = (\Lambda^{-1})^\alpha_\mu \partial_\alpha \phi(x)
\end{equation}

Therefore $\partial_\mu \phi(x)$ is a rank (0,1) tensor.

However, for a contravariant vector $V^\rho$ with $V'^\rho = \Lambda^\rho_\beta V^\beta$:
\begin{equation}
\partial_\mu(V'^\rho) = (\Lambda^{-1})^\alpha_\mu \partial_\alpha(\Lambda^\rho_\beta V^\beta)
\end{equation}
\begin{equation}
= (\Lambda^{-1})^\alpha_\mu \Lambda^\rho_\beta (\partial_\alpha V^\beta) + (\Lambda^{-1})^\alpha_\mu (\partial_\alpha \Lambda^\rho_\beta) V^\beta
\end{equation}

Since $\Lambda^\rho_\beta = \Lambda^\rho_\beta(x)$ in general coordinate transformations, the second term does not vanish, so $\partial_\mu V^\rho$ does \textbf{not} transform as a rank (1,1) tensor.

We need a new derivative operator $\nabla_\mu$ such that $\nabla_\mu V^\rho$ transforms properly as a tensor.

\subsection{Definition of Covariant Derivative}

We demand that:
\begin{equation}
\nabla'_\nu V'^\rho = (\Lambda^{-1})^\alpha_\nu \Lambda^\rho_\beta (\nabla_\alpha V^\beta)
\end{equation}

In a locally inertial frame, the covariant derivative reduces to the ordinary derivative:
\begin{equation}
\nabla_\mu = \partial_\mu \quad \text{(in locally inertial frame)}
\end{equation}

Through careful analysis of the transformation properties, we define:

\begin{equation}
\boxed{\nabla_\alpha V^\rho = \partial_\alpha V^\rho + \Gamma^\rho_{\alpha\beta} V^\beta}
\end{equation}

where $\Gamma^\rho_{\alpha\beta}$ are the \textbf{Christoffel symbols} or \textbf{connection coefficients}.

For covariant vectors:
\begin{equation}
\boxed{\nabla_\alpha V_\beta = \partial_\alpha V_\beta - \Gamma^\rho_{\alpha\beta} V_\rho}
\end{equation}

For scalars:
\begin{equation}
\nabla_\mu \phi = \partial_\mu \phi
\end{equation}

For a mixed tensor $T^\alpha_\beta$:
\begin{equation}
\nabla_\mu T^\alpha_\beta = \partial_\mu T^\alpha_\beta + \Gamma^\alpha_{\mu\rho} T^\rho_\beta - \Gamma^\rho_{\mu\beta} T^\alpha_\rho
\end{equation}

\textbf{General Rule:} The covariant derivative transforms a tensor of rank $(p,q)$ to a tensor of rank $(p, q+1)$.

\subsection{Metric Compatibility}

A fundamental requirement is that the metric is covariantly conserved:
\begin{equation}
\boxed{\nabla_\mu g_{\alpha\beta} = 0}
\end{equation}

Substituting the definition:
\begin{equation}
\nabla_\mu g_{\alpha\beta} = \partial_\mu g_{\alpha\beta} - \Gamma^\rho_{\mu\alpha} g_{\rho\beta} - \Gamma^\rho_{\mu\beta} g_{\alpha\rho} = 0
\end{equation}

By cyclic permutation of indices and solving, we obtain:

\begin{equation}
\boxed{\Gamma^\rho_{\mu\alpha} = \frac{1}{2} g^{\rho\sigma} \left[\partial_\mu g_{\sigma\alpha} + \partial_\alpha g_{\mu\sigma} - \partial_\sigma g_{\mu\alpha}\right]}
\end{equation}

This is the fundamental formula for Christoffel symbols in terms of metric derivatives.

\subsection{Useful Identities}

\begin{align}
\Gamma^\alpha_{\mu\alpha} &= \frac{1}{\sqrt{-g}} \partial_\mu (\sqrt{-g}) \\
\nabla_\mu V^\mu &= \frac{1}{\sqrt{-g}} \partial_\mu (\sqrt{-g} V^\mu) \quad \text{(covariant divergence)}
\end{align}

\subsection{Leibniz Rule}

The covariant derivative satisfies the Leibniz (product) rule:
\begin{align}
\nabla_\mu (\phi V^\nu) &= (\nabla_\mu \phi) V^\nu + \phi \nabla_\mu V^\nu \\
\nabla_\mu (V^\alpha W_\beta) &= (\nabla_\mu V^\alpha) W_\beta + V^\alpha (\nabla_\mu W_\beta)
\end{align}

\section{Curvature: The Riemann Tensor}

\subsection{Commutator of Covariant Derivatives}

For any vector $V^\rho$, consider the commutator:
\begin{equation}
[\nabla_\mu, \nabla_\nu] V^\rho = \nabla_\mu (\nabla_\nu V^\rho) - \nabla_\nu (\nabla_\mu V^\rho)
\end{equation}

After lengthy calculation:
\begin{equation}
[\nabla_\mu, \nabla_\nu] V^\rho = R^\rho_{\sigma\mu\nu} V^\sigma
\end{equation}

where the \textbf{Riemann curvature tensor} is defined as:

\begin{equation}
\boxed{R^\rho_{\sigma\mu\nu} = \partial_\mu \Gamma^\rho_{\nu\sigma} - \partial_\nu \Gamma^\rho_{\mu\sigma} + \Gamma^\rho_{\mu\lambda} \Gamma^\lambda_{\nu\sigma} - \Gamma^\rho_{\nu\lambda} \Gamma^\lambda_{\mu\sigma}}
\end{equation}

This tensor contains at most two derivatives of the metric and measures the curvature of spacetime.

\subsection{Properties of the Riemann Tensor}

\begin{property}[Antisymmetry in last two indices]
\begin{equation}
R^\rho_{\sigma\mu\nu} = -R^\rho_{\sigma\nu\mu}
\end{equation}
\end{property}

\begin{property}[Antisymmetry in first pair]
With all indices lowered: $R_{\rho\sigma\mu\nu} = g_{\rho\lambda} R^\lambda_{\sigma\mu\nu}$
\begin{equation}
R_{\rho\sigma\mu\nu} = -R_{\sigma\rho\mu\nu}
\end{equation}
\end{property}

\textbf{Proof of Property 2:}
Using metric compatibility:
\begin{equation}
0 = [\nabla_\mu, \nabla_\nu] g_{\alpha\beta} = [\nabla_\mu, \nabla_\nu] (V_\alpha W_\beta)
\end{equation}

This leads to:
\begin{equation}
R_{\alpha\beta\mu\nu} + R_{\beta\alpha\mu\nu} = 0
\end{equation}

For covariant vectors:
\begin{equation}
[\nabla_\mu, \nabla_\nu] V_\rho = -R_{\rho\sigma\mu\nu} V^\sigma
\end{equation}

\textbf{Note:} For Minkowski spacetime, $\Gamma^\mu_{\alpha\beta} = 0$ everywhere, so $R^\rho_{\sigma\mu\nu} = 0$.

\subsection{Ricci Tensor and Ricci Scalar}

\begin{definition}[Ricci Tensor]
The Ricci tensor is obtained by contracting the first and third indices:
\begin{equation}
R_{\mu\nu} = R^\rho_{\mu\rho\nu} = g^{\alpha\beta} R_{\alpha\mu\beta\nu}
\end{equation}
\end{definition}

\begin{definition}[Ricci Scalar]
The Ricci scalar (or scalar curvature) is:
\begin{equation}
R = g^{\mu\nu} R_{\mu\nu}
\end{equation}
\end{definition}

The Ricci scalar is a scalar quantity with at most two derivatives of the metric. It plays a central role in the Einstein-Hilbert action.

\section{The Einstein-Hilbert Action}

\subsection{Requirements for Gravitational Action}

The action for the metric $g_{\mu\nu}$ must satisfy:

\begin{enumerate}
\item It should be a scalar under general coordinate transformations
\item It should contain at most 2 derivatives of $g_{\mu\nu}$ (to match Maxwell's theory)
\item It should reduce to Newtonian gravity in the appropriate limit
\end{enumerate}

General form:
\begin{equation}
S_G = \int d^4x \sqrt{-g} \, f(g_{\mu\nu}, \partial_\alpha g_{\mu\nu}, \partial_\alpha\partial_\beta g_{\mu\nu})
\end{equation}

The function $f$ must be a scalar. The only scalar with at most two derivatives is the Ricci scalar $R$ (and constants).

\subsection{Dimensional Analysis}

The action has dimensions:
\begin{equation}
[S] = [ML^2T^{-1}]
\end{equation}

The metric is dimensionless: $[g_{\mu\nu}] = 1$

Newton's constant: $[G] = M^{-1}L^3T^{-2}$

The Ricci scalar: $[R] = L^{-2}$

For the action to have correct dimensions:
\begin{equation}
[\alpha] = \frac{[S]}{L^4 [R]} = \frac{ML^2T^{-1}}{L^4 \cdot L^{-2}} = \frac{ML^2T^{-1}}{L^2} = MT^{-1} = \frac{c^3}{G}
\end{equation}

\subsection{The Einstein-Hilbert Action}

Including a cosmological constant term:

\begin{equation}
\boxed{S_{EH} = \frac{c^3}{16\pi G} \int d^4x \sqrt{-g} \, (R - 2\lambda)}
\end{equation}

where:
\begin{itemize}
\item $\lambda$ is the \textbf{cosmological constant}
\item The factor $16\pi$ is chosen to reproduce Newtonian gravity in the weak-field limit
\item This action contains at most 2 derivatives of the metric
\end{itemize}

Setting $c = 1$ (natural units), we often write:
\begin{equation}
S_{EH} = \frac{1}{16\pi G} \int d^4x \sqrt{-g} \, R
\end{equation}

Or with convenient normalization:
\begin{equation}
S_{EH} = \frac{1}{2} \int d^4x \sqrt{-g} \, R
\end{equation}

\section{Scalar Field Theory in Curved Spacetime}

\subsection{Action for Minimally Coupled Scalar Field}

\subsubsection{Flat Spacetime Action}

In flat Minkowski spacetime, the action for a real scalar field is:
\begin{equation}
S_\phi^{\text{flat}} = \int d^4x \left[ \frac{1}{2} \partial^\mu \phi \partial_\mu \phi - V(\phi) \right]
\end{equation}

Requirements for the Lagrangian density:
\begin{itemize}
\item $\mathcal{L}$ must be a scalar
\item At most two derivatives in equations of motion
\item No explicit coordinate dependence (translational invariance)
\end{itemize}

General form:
\begin{equation}
\mathcal{L} = A(\phi)\partial_\mu \phi \partial^\mu \phi - U(\phi)
\end{equation}

By field redefinition $\chi(\phi) = \int^\phi d\varphi \sqrt{2A(\varphi)}$, this can always be brought to canonical form:
\begin{equation}
\mathcal{L} = \frac{1}{2} \partial_\mu \chi \partial^\mu \chi - V(\chi)
\end{equation}

\subsubsection{Curved Spacetime Action}

To promote the scalar field action to curved spacetime (minimal coupling):

\begin{itemize}
\item $\partial^\mu \phi \partial_\mu \phi \to g^{\mu\nu} \partial_\mu \phi \partial_\nu \phi$
\item $d^4x \to \sqrt{-g} \, d^4x$
\item $\partial_\mu \to \nabla_\mu$ (though for scalars this makes no difference)
\end{itemize}

The curved spacetime action becomes:
\begin{equation}
\boxed{S_\phi = \int d^4x \sqrt{-g} \left[ \frac{1}{2} g^{\mu\nu} \partial_\mu \phi \partial_\nu \phi - V(\phi) \right]}
\end{equation}

\subsection{Equations of Motion}

Varying with respect to $\phi$ and setting $\delta S = 0$:

\begin{equation}
\delta S_\phi = \int d^4x \sqrt{-g} \left[ g^{\mu\nu} (\partial_\mu \phi) \partial_\nu (\delta\phi) - V_{,\phi} \delta\phi \right]
\end{equation}

Integrating by parts and requiring $\delta S_\phi = 0$ for arbitrary $\delta\phi$:

\begin{equation}
\boxed{\Box_g \phi + V_{,\phi} = 0}
\end{equation}

where the curved spacetime d'Alembertian operator is:
\begin{equation}
\Box_g = \frac{1}{\sqrt{-g}} \partial_\mu \left(\sqrt{-g} \, g^{\mu\nu} \partial_\nu\right)
\end{equation}

This is the \textbf{covariant Klein-Gordon equation} in curved spacetime.

\subsection{Full Action: Gravity + Matter}

The complete action coupling gravity to a scalar field:

\begin{equation}
\boxed{S = S_{EH} + S_\phi = \int d^4x \sqrt{-g} \left[ \frac{1}{2} R + \frac{1}{2} g^{\mu\nu} \partial_\mu \phi \partial_\nu \phi - V(\phi) \right]}
\end{equation}

\section{Scalar Field Cosmology}

\subsection{FRW Metric}

The Friedmann-Robertson-Walker (FRW) metric for a homogeneous and isotropic universe:

\begin{equation}
ds^2 = -dt^2 + a^2(t) \left[ dr^2 + r^2 d\Omega^2 \right] = -dt^2 + a^2(t) d\vec{x}^2
\end{equation}

The metric components:
\begin{equation}
g_{\mu\nu} = \text{diag}(-1, a^2, a^2, a^2)
\end{equation}

The determinant:
\begin{equation}
g = \det(g_{\mu\nu}) = -a^6 \quad \Rightarrow \quad \sqrt{-g} = a^3
\end{equation}

\subsection{Friedmann-Klein-Gordon Equation}

For a homogeneous scalar field $\phi = \phi(t)$:
\begin{equation}
g^{\mu\nu} \partial_\mu \phi \partial_\nu \phi = -\dot{\phi}^2
\end{equation}

Computing the d'Alembertian:
\begin{equation}
\Box_g \phi = \frac{1}{a^3} \partial_t (a^3 \times (-\dot{\phi})) = -\ddot{\phi} - 3\frac{\dot{a}}{a}\dot{\phi}
\end{equation}

The Klein-Gordon equation $\Box_g \phi + V_{,\phi} = 0$ becomes:

\begin{equation}
\boxed{\ddot{\phi} + 3H\dot{\phi} + V_{,\phi} = 0}
\end{equation}

where $H = \dot{a}/a$ is the Hubble parameter.

\textbf{Physical interpretation:}
\begin{itemize}
\item $\ddot{\phi}$ : acceleration of the field
\item $3H\dot{\phi}$ : Hubble friction (damping due to cosmological expansion)
\item $V_{,\phi}$ : restoring force from potential gradient
\end{itemize}

\subsection{Stress-Energy Tensor for Scalar Field}

The stress-energy tensor is obtained by varying the matter action with respect to the metric:
\begin{equation}
T_{\mu\nu} = \frac{-2}{\sqrt{-g}} \frac{\delta S_\phi}{\delta g^{\mu\nu}} = \partial_\mu \phi \partial_\nu \phi - g_{\mu\nu} \left( \frac{1}{2} \partial^\sigma \phi \partial_\sigma \phi - V(\phi) \right)
\end{equation}

For a homogeneous field in FRW cosmology, this takes the form of a perfect fluid:

\textbf{Energy density:}
\begin{equation}
\rho_\phi = \frac{1}{2} \dot{\phi}^2 + V(\phi)
\end{equation}

\textbf{Pressure:}
\begin{equation}
P_\phi = \frac{1}{2} \dot{\phi}^2 - V(\phi)
\end{equation}

\textbf{Equation of state parameter:}
\begin{equation}
w_\phi = \frac{P_\phi}{\rho_\phi} = \frac{\frac{1}{2}\dot{\phi}^2 - V}{\frac{1}{2}\dot{\phi}^2 + V}
\end{equation}

\subsection{Inflationary Scenarios}

When the potential energy dominates ($V \gg \frac{1}{2}\dot{\phi}^2$):
\begin{itemize}
\item $w_\phi < 0$ : negative pressure
\item $w_\phi < -\frac{1}{3}$ : accelerated expansion
\end{itemize}

This provides a mechanism for cosmological inflation.

\subsection{Friedmann Equations}

The dynamics of the scalar field coupled to gravity are governed by:

\begin{align}
\ddot{\phi} + 3H\dot{\phi} + V_{,\phi} &= 0 \quad \text{(Klein-Gordon)} \\
H^2 &= \frac{1}{3} \left( \frac{1}{2}\dot{\phi}^2 + V(\phi) \right) \quad \text{(Friedmann, assuming $k=0$)}
\end{align}

These equations describe the evolution of a homogeneous scalar field in an expanding FRW universe (flat spatial sections, $k=0$).

\section{Maxwell's Equations in Curved Spacetime}

\subsection{Covariant Form}

The Maxwell equations in curved spacetime can be written using the covariant divergence:

\begin{equation}
\boxed{\nabla_\mu F^{\mu\nu} = \mu_0 J^\nu}
\end{equation}

where the covariant divergence is:
\begin{equation}
\nabla_\mu F^{\mu\nu} = \frac{1}{\sqrt{-g}} \partial_\mu \left(\sqrt{-g} \, F^{\mu\nu}\right)
\end{equation}

Explicitly:
\begin{equation}
\partial_\mu \left(\sqrt{-g} \, F^{\mu\nu}\right) = \sqrt{-g} \, \mu_0 J^\nu
\end{equation}

In flat spacetime ($\sqrt{-g} = 1$), this reduces to:
\begin{equation}
\partial_\mu F^{\mu\nu} = \mu_0 J^\nu
\end{equation}

\subsection{Important Note on Field Tensor}

The electromagnetic field tensor is defined as:
\begin{equation}
F_{\mu\nu} = \nabla_\mu A_\nu - \nabla_\nu A_\mu = \partial_\mu A_\nu - \partial_\nu A_\mu
\end{equation}

The Christoffel symbols cancel in the antisymmetric combination, so the covariant and ordinary derivatives give the same result for $F_{\mu\nu}$.

\section{Summary of Key Results}

\subsection{Special Relativity}
\begin{itemize}
\item Lorentz transformations with $\gamma = (1-v^2/c^2)^{-1/2}$
\item Four-vectors: position, velocity, momentum, force
\item Electromagnetic field tensor $F^{\mu\nu}$ and covariant Maxwell equations
\item Gauge invariance and electromagnetic action
\end{itemize}

\subsection{General Relativity}
\begin{itemize}
\item Invariant volume element: $dV = \sqrt{-g} \, d^4x$
\item Covariant derivative with Christoffel symbols
\item Metric compatibility: $\nabla_\mu g_{\alpha\beta} = 0$
\item Riemann curvature tensor measures spacetime curvature
\item Einstein-Hilbert action: $S_{EH} = \frac{c^3}{16\pi G} \int d^4x \sqrt{-g} (R - 2\lambda)$
\end{itemize}

\subsection{Scalar Field Cosmology}
\begin{itemize}
\item Minimally coupled scalar field action in curved spacetime
\item Covariant Klein-Gordon equation
\item Friedmann-Klein-Gordon equation for homogeneous fields
\item Scalar field as perfect fluid with equation of state parameter $w_\phi$
\item Mechanism for inflation when potential dominates
\end{itemize}

\end{document}